\section{Exact bounds near the two standard constructions}\label{sec:exact}

The stability statement leaves an error of order $o(n^k)$. To prove \cref{thm:matching}, we need exact upper bounds for hypergraphs close to each of the two standard constructions. We first record a consequence of the classical theorem of Bollob\'as, Daykin, and Erd\H{o}s~\cite[Theorem 1]{BDE1976}, which gives the bound and equality case under the weaker assumption $n>2k^3b$ for $b\ge1$. The self-contained induction below follows the high-degree-vertex argument of Erd\H{o}s~\cite[pp.~94--95]{Erdos1965}.

\begin{lemma}\label{lem:sparse-matching}
Let $k\ge2$ and $b\ge0$ be integers. If $H$ is an $n$-vertex $k$-graph, $\nu(H)\le b$, and $n\ge4k^3b$, then $|H|\le M(n,b)$. Equality holds only for a complete $b$-cover.
\end{lemma}
\begin{proof}
We induct on $b$. If $b=0$, then $H$ has no edge. Suppose $b\ge1$. We may assume $|H|\ge M(n,b)$, since a smaller family already satisfies the strict bound. The vertices of a maximal matching meet every edge, and there are at most $kb$ such vertices. Some vertex $v$ among them has degree
\begin{align}
d_H(v)\ge\frac{|H|}{kb}
\ge\frac1k\binom{n-b}{k-1}
>kb\binom{n-2}{k-2}.
\label{eq:sparse-large-degree}
\end{align}
For the middle inequality, write $M(n,b)=\sum_{j\in[b]}\binom{n-j}{k-1}\ge b\binom{n-b}{k-1}$. For the strict inequality, the product formula for binomial coefficients gives
\begin{align*}
\frac{\binom{n-b}{k-1}}{\binom{n-1}{k-1}}
=\prod_{0\le h\le k-2}\left(1-\frac{b-1}{n-1-h}\right)
\ge1-\frac{(k-1)(b-1)}{n-k+1}>\frac12.
\end{align*}
The assumption on $n$ also implies $n-1>2k^2(k-1)b$. Since $\binom{n-1}{k-1}=\frac{n-1}{k-1}\binom{n-2}{k-2}$, these two bounds prove \cref{eq:sparse-large-degree}.

If $H-v$ had a matching of size $b$, the number of edges through $v$ meeting its $kb$ vertices would be at most $kb\binom{n-2}{k-2}$. By \cref{eq:sparse-large-degree}, an edge through $v$ avoids the matching, contradicting $\nu(H)\le b$. Thus $\nu(H-v)\le b-1$. Since $n-1\ge4k^3(b-1)$, induction yields
\begin{align*}
|H|\le\binom{n-1}{k-1}+M(n-1,b-1)=M(n,b).
\end{align*}
Equality requires every $k$-set through $v$ to be present and $H-v$ to be a complete $(b-1)$-cover. Their union is a complete $b$-cover.
\end{proof}

\subsection{Hypergraphs close to a complete cover}

For a vertex in the center set of a complete cover, the link on the outside vertices is complete. If only a few of these link edges are missing, that vertex can still extend a prescribed matching. The following lemma is the form with a bound only on missing cover edges in Chi, Sun, and Wang~\cite[Lemma 8.6]{CSW2026}. Its argument follows Cao, Liu, and Zhang~\cite[Lemma 5.2]{CLZ2026}, whose statement assumes closeness in symmetric difference.
\begin{lemma}\label{lem:cover-exact}
Fix $k\ge3$, $\gamma>0$, and $K\ge k+\gamma$. There are $\eta>0$ and $s_0$ such that the following holds. Suppose $s\ge s_0$, $(k+\gamma)s\le n\le Ks$, and $H$ is an $n$-vertex $k$-graph with $\nu(H)\le s$. If an $s$-set $S$ satisfies
\begin{align*}
|\mathcal A(V(H),S)\setminus H|\le\eta n^k,
\end{align*}
then $|H|\le M(n,s)$. Equality holds only for a complete $s$-cover, whose center set need not be the original set $S$.
\end{lemma}
\begin{proof}
Set $W=V(H)\setminus S$. Choose
\begin{align*}
\beta\coloneqq\min\{\gamma/4,(k-1)/(8k^3),1/4\},
\qquad
c_0\coloneqq\frac{(\gamma/4)^{k-1}}{(k-1)!},
\end{align*}
and take $\eta>0$ with $\eta K^k/c_0\le\beta$. For sufficiently large $s$, put
\begin{align*}
D_0&\coloneqq\binom{\lfloor\gamma s/2\rfloor}{k-1},\\
B&\coloneqq\left\{v\in S\colon
\left|\{A\in\tbinom W{k-1}\colon A\cup\{v\}\notin H\}\right|\ge D_0\right\}.
\end{align*}
The sets counted for different vertices $v\in S$ are different missing edges, because each contains exactly one vertex of $S$. Also $D_0\ge c_0s^{k-1}$ for large $s$. Hence, writing $b=|B|$, we have
\begin{align*}
bD_0\le\eta n^k\le\eta K^ks^k,
\qquad b\le\beta s.
\end{align*}

Let $U=W\cup B$. We claim that $\nu(H[U])\le b$. Suppose instead that $H[U]$ contains a matching of size $b+1$. Extend it in turn through every vertex of $S\setminus B$, using $k-1$ unused vertices of $W$ for each new edge. Before each step, the matching has at most $s$ edges. The original $b+1$ edges use at most $k(b+1)$ vertices of $W$, and at most $s-b-1$ previously added edges use another $(k-1)(s-b-1)$. The number of available vertices of $W$ is therefore at least
\begin{align*}
(n-s)-k(b+1)-(k-1)(s-b-1)
&=n-ks-b-1\\
&\ge\gamma s-\beta s-1\ge\lfloor\gamma s/2\rfloor.
\end{align*}
There are at least $D_0$ possible $(k-1)$-sets among them. The current center is not in $B$, so fewer than $D_0$ of these sets fail to form an edge with it. An extension is therefore possible at every step. We obtain $b+1+(s-b)=s+1$ disjoint edges, a contradiction.

Since $|U|\ge n-s\ge(k-1)s$ and $b\le\beta s$, we have $|U|\ge4k^3b$. Applying \cref{lem:sparse-matching} to $H[U]$ gives
\begin{align*}
|H|&\le\binom nk-\binom{|U|}k+M(|U|,b)\\
&=\binom nk-\binom{|U|-b}k
=\binom nk-\binom{n-s}k.
\end{align*}
If equality holds, every $k$-set meeting $S\setminus B$ belongs to $H$, and $H[U]$ is a complete $b$-cover. If $T\subseteq U$ is its set of $b$ centers, then $H=\mathcal A(V(H),(S\setminus B)\cup T)$.
\end{proof}

\subsection{Hypergraphs containing a large clique}

We next prove the exact bound when a genuine large clique is available. This is a variant of the large-clique step in Chi, Sun, and Wang~\cite[Lemma 8.7]{CSW2026}, with a different clique-size threshold and a proof using random disjoint blocks and a bipartite matching bound. Its vertices can complete a partial matching. This turns the matching restriction into a bound on the number of outside vertices covered by disjoint edges.
\begin{lemma}\label{lem:local-clique}
Fix $k\ge3$. For all sufficiently large $s$, let $k(s+1)\le n\le(k+1/2)s$, and let $H$ be an $n$-vertex $k$-graph with $\nu(H)\le s$. If $H$ contains a complete $k$-graph on
\begin{align*}
u\coloneqq\lfloor(k-1/6)s\rfloor
\end{align*}
vertices, then $|H|\le C(s)$. Equality holds only when $H$ is a complete $k$-graph on $ks+k-1$ vertices, together with isolated vertices.
\end{lemma}
\begin{proof}
Let $Q$ be the vertex set of the given clique, and put
\begin{align*}
W\coloneqq V(H)\setminus Q,\qquad
N\coloneqq|W|,\qquad
d\coloneqq ks+k-1-u,\qquad
C\coloneqq\binom{u}{k-1}.
\end{align*}
For fixed $k$ and sufficiently large $s$, the assumptions imply
\begin{align}
&1\le d<N,\qquad N\to\infty,\qquad u\ge\frac32(k-1)N,\nonumber\\
&u\ge(k-1)(d+1),\qquad u\ge N(k-1)+1.
\label{eq:clique-size-conditions}
\end{align}
For example, $N\le(2/3)s+1$ while $u\ge(k-1/6)s-1$, so the third inequality has a positive linear difference between its sides for every fixed $k\ge3$. The other inequalities follow from $n\ge ks+k$, $d=s/6+O_k(1)$, and the same estimates.

Every matching consisting of edges that meet $W$ covers at most $d$ vertices of $W$. To see this, let such a matching have $h$ edges and cover $w>d$ outside vertices. We already have $h\le s$. Its edges use $kh-w$ vertices of $Q$, leaving at least
\begin{align*}
u-kh+w\ge u-kh+d+1=k(s+1-h)
\end{align*}
unused vertices of the clique. Partitioning this many vertices into $s+1-h$ $k$-sets extends the matching to size $s+1$, a contradiction.

Let $E_j$ be the number of edges of $H$ containing exactly $j$ vertices of $W$, for $0\le j\le k$. We first bound $E_1$. Choose uniformly an ordered collection of $N$ pairwise disjoint $(k-1)$-sets in $Q$. Form a bipartite graph $G$ whose two vertex classes are $W$ and these $N$ sets, joining $v$ to $A$ when $A\cup\{v\}\in H$. A matching in $G$ gives disjoint hyperedges using one distinct outside vertex each, so $\nu(G)\le d$. Write $\nu(G)=d-e$, where $e\ge0$ depends on the sample. By K\H{o}nig's theorem, $G$ has a vertex cover of size $d-e$, and hence at most $N(d-e)$ edges. Each sampled block is a uniform $(k-1)$-set of $Q$, so $\E|G|=NE_1/C$. It follows that
\begin{align}
D\coloneqq dC-E_1\ge0,\qquad
\E e\le d-\frac{\E|G|}{N}=\frac DC.
\label{eq:singleton-deficit}
\end{align}

Fix $2\le j\le k$ and set $b=\lfloor N/j\rfloor$. In addition to the first $N$ blocks, sample $b$ pairwise disjoint $(k-j)$-sets $B_1,\ldots,B_b$ in $Q$, disjoint from those blocks. The sampling is possible because
\begin{align*}
(k-1)N+(k-j)b\le\tfrac32(k-1)N\le u.
\end{align*}
The joint distribution is uniform among all such disjoint collections. In particular, each $B_\ell$ is uniform among the $(k-j)$-sets of $Q$; independence between blocks is not needed. When $j=k$, the sets $B_\ell$ are empty.

Condition on all these blocks, and choose a maximum matching in $G$. Let $T\subseteq W$ be its set of matched vertices, so $|T|=d-e$. Independently choose $b$ uniformly random disjoint $j$-sets $I_1,\ldots,I_b$ of $W$. Insert all tested edges $I_\ell\cup B_\ell$ that belong to $H$, deleting the singleton-outside matching edges whose outside vertices lie in an inserted $I_\ell$. The new edges have disjoint inside blocks, so no further conflict occurs. The number of additional covered outside vertices is
\begin{align*}
\sum_{\ell\in[b]}|I_\ell\setminus T|\,\1_{\{I_\ell\cup B_\ell\in H\}}\le e,
\end{align*}
by the outside-vertex bound proved above. Each successful test with $I_\ell\nsubseteq T$ contributes at least one. Each test has probability at most $\binom dj/\binom Nj$ of lying entirely in $T$. Averaging first over the $I_\ell$ and then over the inside blocks, and using \cref{eq:singleton-deficit}, yields
\begin{align*}
\frac{bE_j}{\binom Nj\binom{u}{k-j}}
\le\frac{b\binom dj}{\binom Nj}+\E e
\le\frac{b\binom dj}{\binom Nj}+\frac DC.
\end{align*}
Equivalently,
\begin{align}
E_j\le\binom dj\binom{u}{k-j}
+\frac{\binom Nj}{\lfloor N/j\rfloor}\binom{u}{k-j}\frac DC.
\label{eq:multiple-outside-count}
\end{align}

We now sum these estimates. Since $Q$ is complete, $E_0=\binom uk$, while $E_1=dC-D$. Set
\begin{align*}
\Theta\coloneqq\frac1C\sum_{2\le j\le k}
\frac{\binom Nj}{\lfloor N/j\rfloor}\binom{u}{k-j}.
\end{align*}
Vandermonde's identity and \cref{eq:multiple-outside-count} give
\begin{align}
|H|\le\binom{u+d}k-(1-\Theta)D.
\label{eq:local-clique-deficit}
\end{align}
It remains to check that $\Theta<1$. For $2\le j\le k$, we have $\lfloor N/j\rfloor\ge(N-j+1)/j\ge(N-k+1)/j$. Hence
\begin{align*}
\Theta
&\le\frac{N}{N-k+1}\left(\frac{\binom{u+N-1}{k-1}}{\binom{u}{k-1}}-1\right),\\
\frac{\binom{u+N-1}{k-1}}{\binom{u}{k-1}}
&=\prod_{0\le h\le k-2}\left(1+\frac{N-1}{u-h}\right)
\le\exp\left(\frac{(k-1)(N-1)}{u-k+2}\right)<e^{2/3}.
\end{align*}
The last strict inequality follows from $u\ge\frac32(k-1)N$. Since $N\to\infty$ and $e^{2/3}-1<1$, we indeed have $\Theta<1$ for sufficiently large $s$. Now $u+d=ks+k-1$, so \cref{eq:local-clique-deficit} proves the required upper bound.

Suppose equality holds. Then $D=0$. In the first sampling, every graph $G$ has at most $Nd$ edges, and their average is exactly $Nd$ by \cref{eq:singleton-deficit}. Thus every sample has exactly $Nd$ edges. A minimum vertex cover consisting of $p$ vertices on one side and $q$ on the other has $p+q=d$ and covers at most $N(p+q)-pq$ edges. Equality forces $pq=0$. Consequently every sampled graph consists either of $d$ complete rows or of $d$ complete columns.

For a $(k-1)$-set $A\subseteq Q$, define
\begin{align*}
I(A)\coloneqq\{v\in W\colon A\cup\{v\}\in H\}.
\end{align*}
If some $I(A)$ is nonempty and proper, every sample containing $A$ must have $d$ complete rows, so $I(A)$ is a fixed $d$-set $T$. Any $(k-1)$-set disjoint from $A$ can be included with $A$ in a sample, and therefore has the same set $T$ of neighbors. For an arbitrary $(k-1)$-set $A'$, choose a $(k-1)$-set disjoint from $A\cup A'$, which is possible by \cref{eq:clique-size-conditions}. Applying the preceding observation twice gives $I(A')=T$.

If no $I(A)$ is nonempty and proper, each is either empty or all of $W$. Two $(k-1)$-sets that differ in one vertex can be completed by the same $N-1$ disjoint blocks, because $u\ge N(k-1)+1$. Every sample has exactly $d$ full columns, so the two sets must have the same type. Any two $(k-1)$-sets can be connected by one-vertex changes, forcing all $I(A)$ to be empty or all to equal $W$. This would give zero or $N$ full columns, contrary to $0<d<N$.

It follows that exactly $d$ outside vertices, forming a set $T$, have complete links on $Q$, and every other outside vertex has empty singleton link. Suppose an edge $E$ nevertheless meets $W\setminus T$. For each $v\in T\setminus E$, choose a fresh $(k-1)$-set in $Q\setminus E$ and use its complete link with $v$. The choices can be made disjoint because $u\ge(k-1)(d+1)$ and $|E\cap Q|\le k-1$. Together with $E$, these edges form a matching covering
\begin{align*}
|E\cap W|+|T\setminus E|=d+|E\cap(W\setminus T)|>d
\end{align*}
outside vertices, a contradiction. Therefore every edge lies in $Q\cup T$. Equality of the edge count forces the whole clique on this $(ks+k-1)$-set.
\end{proof}

\subsection{Shifting and the conclusion}

To use \cref{lem:local-clique}, we must obtain a genuine clique from a family that is only close to one. We use shifts only inside the proposed clique vertex set. For two such vertices $i<j$, replace an edge $E$ containing $j$ and not $i$ by $(E\setminus\{j\})\cup\{i\}$ if that latter edge is absent; otherwise leave $E$ unchanged. This operation preserves the number of edges and the number of edges contained in the proposed clique.

It also cannot increase the matching number. Indeed, a matching in the shifted family contains at most one new edge, since every new edge contains $i$. If the matching has a new edge $E'$, replace it by the original edge obtained by exchanging $i$ for $j$. This is already a matching unless another edge $F$ uses $j$. In that case, the survival of $F$ under the shift means that $(F\setminus\{j\})\cup\{i\}$ was also an original edge. Replacing the two edges simultaneously gives an original matching of the same size.

The next lemma follows the shifting and equality argument of Chi, Sun, and Wang~\cite[Theorem 8.8]{CSW2026}, using \cref{lem:local-clique} for the large-clique step. Its assumption involves only missing edges inside the proposed clique.
\begin{lemma}\label{lem:clique-exact}
Fix $k\ge3$. There are $\eta>0$ and $s_0$ such that the following holds. Let $s\ge s_0$, $k(s+1)\le n\le(k+1/2)s$, and let $H$ be an $n$-vertex $k$-graph with $\nu(H)\le s$. If a set $Q_0$ of size $ks+k-1$ satisfies
\begin{align*}
\left|\binom{Q_0}k\setminus H\right|\le\eta s^k,
\end{align*}
then $|H|\le C(s)$, with equality only for a complete $(ks+k-1)$-vertex clique.
\end{lemma}
\begin{proof}
Put $m=ks+k-1$ and $u=\lfloor(k-1/6)s\rfloor$, and label $Q_0$ by $[m]$. Apply shifts within $Q_0$ until none changes the family. A nontrivial shift decreases the sum, over all edges, of the labels of their vertices in $Q_0$, so the process terminates. Let $H'$ be the resulting family. It has the same number of edges and the same number of missing edges inside $Q_0$ as $H$, and $\nu(H')\le s$.

If $H'$ misses a $k$-set contained in $[u]$, then it misses every $k$-set in $Q_0\setminus[u]$. Indeed, each of those sets lies above the missing set in the componentwise order, and successive shifts inside $Q_0$ would otherwise imply that the missing set is present. Thus $H'$ would miss at least $\binom{m-u}k$ edges of $Q_0$. Since $m-u=s/6+O_k(1)$, we have $\binom{m-u}k\ge s^k/(12^kk!)$ for sufficiently large $s$. Taking $\eta<1/(12^kk!)$ excludes this possibility. Consequently $[u]$ spans a complete $k$-graph in $H'$, and \cref{lem:local-clique} gives $|H|=|H'|\le C(s)$.

If equality holds, \cref{lem:local-clique} says that $H'$ is an $m$-vertex clique. Consider the last nontrivial shift in the sequence, with shifted vertices $i,j$. Just before that shift, every edge of this clique avoiding $i,j$ was already present. Thus the preceding family contains a clique on at least $m-2\ge u$ vertices. It has matching number at most $s$ and size $C(s)$, so the equality case of \cref{lem:local-clique} implies that it too is an $m$-vertex clique. Repeat this argument backwards through the finite shift sequence. The original family $H$ is therefore an $m$-vertex clique.
\end{proof}

We can now combine stability with the two local bounds. The argument also explains why the threshold on $s$ is uniform in $n$.
\begin{proof}[Proof of \cref{thm:matching}]
Suppose the conclusion fails for arbitrarily large $s$. Then there is a sequence of nonstandard families $H_i$ with $s_i\to\infty$, $n_i\ge k(s_i+1)$, $\nu(H_i)\le s_i$, and $|H_i|\ge D(n_i,s_i)$. Here a family is nonstandard if it violates the numerical bound or attains the bound without having either stated extremal form.

By \cref{lem:sparse-matching}, none of these counterexamples has $n_i\ge4k^3s_i$: in that range the lemma gives $|H_i|\le M(n_i,s_i)\le D(n_i,s_i)$, and equality forces a standard cover. We may pass to a subsequence such that $n_i/s_i\to R\in[k,4k^3]$. Applying \cref{prop:matching-stability} with errors tending to zero, and then passing to one of its two alternatives, makes $H_i$ either $o(n_i^k)$-close to a complete $s_i$-cover or $o(n_i^k)$-close to a complete $(ks_i+k-1)$-vertex clique.

In the cover case, $R=k$ is impossible: the cover density would tend to $1-(1-1/k)^k<1$, whereas $D(n_i,s_i)/\binom{n_i}k$ would tend to one. For $R>k$, comparison of limiting edge densities gives $R\ge\rho_k>k$, as explained after \cref{prop:matching-stability}. Choose $0<\gamma<R-k$ and take $K=4k^3$. Eventually $(k+\gamma)s_i\le n_i\le Ks_i$, and the number of missing cover edges is smaller than the error allowed by \cref{lem:cover-exact}. Hence
\begin{align*}
D(n_i,s_i)\le|H_i|\le M(n_i,s_i)\le D(n_i,s_i),
\end{align*}
and equality forces a complete cover, contrary to the choice of $H_i$.

In the clique case, comparison of limiting densities gives $R\le\rho_k<k+1/2$, including $R=k$. Thus eventually $n_i\le(k+1/2)s_i$. Since $n_i=O_k(s_i)$, the number of missing clique edges is $o(s_i^k)$. Applying \cref{lem:clique-exact} gives
\begin{align*}
D(n_i,s_i)\le|H_i|\le C(s_i)\le D(n_i,s_i),
\end{align*}
and equality forces a complete clique, again a contradiction. Both possible subsequences have been excluded, proving the theorem.
\end{proof}
